% !PRINTED 64
\weberpage{64}

\begin{center}
{\large\gesperrt{Zweiter Abschnitt.}}\\[0.5em]
{\Large\gesperrt{Determinanten.}}\\[0.4em]
\end{center}

\begin{center}
{\large \S.\ 18.}

\vspace{0.5em}
{\large\gesperrt{Permutationen von \(n\) Elementen.}}
\end{center}

Wir betrachten ein System von \(n\) unterschiedenen Elementen irgend welcher Art, z.~B. die \(n\) Ziffern
\[
1,2,3,\ldots,n,
\]
deren Complex in dieser bestimmten Anordnung wir mit \(\mathfrak A\) bezeichnen wollen. Die Elemente von \(\mathfrak A\) lassen sich auf verschiedene Arten anordnen, z.~B.
\[
2,1,3,\ldots,n.
\]
Der Uebergang von einer Anordnung zu einer anderen heisst eine \gesperrt{Permutation}.

Bezeichnen wir die Anzahl der verschiedenen Anordnungen, die nur von der Anzahl \(n\) der Elemente abhängen kann, mit \(\Pi(n)\), so ergiebt sich zunächst
\[
\Pi(1)=1,\qquad \Pi(2)=2,
\]
und um die Zahl allgemein zu bestimmen, denken wir uns zu \(n-1\) Elementen ein \(n\)tes hinzugefügt. In jeder Anordnung der \(n-1\) Elemente kann nun das \(n\)te Element an \(n\) verschiedene Stellen gesetzt werden, nämlich vor das erste, zwischen das erste und zweite, zwischen das zweite und dritte u.~s.~f., endlich nach dem \((n-1)\)ten, und alle die so entstandenen Anordnungen sind von einander verschieden. Daraus folgt:
\begin{equation}
\Pi(n)=n\Pi(n-1),
\tag{1}
\end{equation}
woraus sich durch vollständige Induction
\begin{equation}
\Pi(n)=1\cdot2\cdot3\cdots n
\tag{2}
\end{equation}
%CONT>>
